\subsection[Description of $q$-TAZRP]{Description of $\boldsymbol{q}$-TAZRP}

Throughout this paper, $q \in (0, 1)$ is a constant. Let $n$ particles be on the integer lattice $\intZ$, and denote the position of the $k$-th particle by $x_k$, or $x_k(t)$ if we specify the time $t \geq 0$. We also use $x_k$ to denote the $k$-th particle itself by abuse of notation. We assume that initially $(x_1(0), x_2(0), \dotsc, x_n(0)) =: X(0) \in \Weylchamber^n$, where $\Weylchamber^n = \{ (i_1, \dotsc, i_n) \in \intZ^n \,|\, i_1 \geq i_2 \geq \dotsb \geq i_n \}$. From the dynamics given below, we know that for all $t > 0$ the order $x_1(t) \geq x_2(t) \geq \dotsb \geq x_n(t)$ is kept, or equivalently, $X(t) \in \Weylchamber^n$.

Each particle is at an integer site at any given time, and several consecutive particles can share a site. Suppose $x_{k - 1}(t) > x_k(t) = x_{k + 1}(t) = \dotsb = x_{k + j - 1}(t) = m > x_{k + j}(t)$, then we say that at time $t$, the particles $x_k, x_{k + 1}, \dotsc, x_{k + j - 1}$ are in a stack at $m$ with height $j$, such that $x_k$ is at the top of the stack, $x_{k + i}$ is below $x_{k + i - 1}$ ($i = 1, \dotsc, j - 1$), and $x_{k + j - 1}$ is at the bottom. See Fig.~\ref{fig:particles}.
\begin{figure}[htb]
 \begin{minipage}[t]{0.46\linewidth}
 \centering
 \includegraphics{particles}
 \caption{Particles in $q$-TAZRP. The $k$-th particle is on the top of a stack at $i - 1$ with height~$3$.}
 \label{fig:particles}
 \end{minipage}
 \hfill
 \begin{minipage}[t]{0.48\linewidth}
 \centering
 \includegraphics{particles_jumped}
 \caption{The $k$-th particle jumps to the bottom of the stack at $i$. Until the $k$-th particle jumps, the $(k + 1)$-th and $(k + 2)$-th particles stay put.}
 \label{fig:particles_jumped}
 \end{minipage}
\end{figure}

On each site $i$, we assign a \emph{conductance} $a_i > 0$. Throughout this paper, we denote $b_i = a_i(1 - q)$ for notational simplicity. The dynamics of the particles is given as follows.
If a particle is in a stack but is not on the top of the stack, then it cannot move; otherwise if the particle is on the top of a stack at $m$ with height $j$, it moves to site $m + 1$ at rate $a_m(1 - q^j) = b_m(1 + q + \dotsb + q^{j - 1})$.
The movement occurs regardless of whether the right neighbour site is occupied or not (hence the model is named ``zero range''). If the right neighbour site is occupied, then $x_k$ moves to the bottom of the stack there, and the height of that stack increases by $1$. See Fig.~\ref{fig:particles_jumped} for illustration.

The $q$-TAZRP is clearly a Markov process. Let $X = (x_1, x_2, \dotsc, x_n) \in \Weylchamber^n$ denote a state of the model, then $P(X; t)$, the probability that the model is at state $X$ at time $t$, satisf\/ies the master equation
\begin{gather} \label{eq:abstract_master_eq}
 \frac{d}{dt} P(X; t) = H P(X; t) = \sum_{Y \in \Weylchamber^n} P(Y; t) (c(Y, X) - c(X, Y)),
\end{gather}
where $H$ is the generator of the Markov process, and $c(X, Y)$ is the transition rate for the Markov process from state $X$ to state $Y$. To express the master equation more concretely, for $X = (x_1, \dotsc, x_n) \in \Weylchamber^n$ we denote
\begin{gather*}
 X^{k, -} = (x_1, \dotsc, x_{k - 1}, x_k - 1, x_{k + 1}, \dotsc, x_n), \qquad k = 1, \dotsc, n,
\end{gather*}
and denote
\begin{gather} \label{eq:various_notations}
 \begin{split}
 & \ell(X) :=  \text{$\#$ of the distinct values attained by $x_1, \dotsc, x_n$}, \\
 & v_i(X) :=  \text{the $i$-th largest value attained by $x_1, \dotsc, x_n$}, \\
& p_k(X) :=  \text{$\#$ of $x_1, \dotsc, x_n$ whose value is $k$}, \\
 & n_i(X) :=  p_{v_i}(X), \quad \text{and} \quad N_i(X) := n_1(X) + \dotsb + n_i(X).
 \end{split}
\end{gather}
For instance, if
\begin{gather} \label{eq:general_X}
 X = (\underbrace{s_1, \dotsc, s_1}_{k_1}, \underbrace{s_2, \dotsc, s_2}_{k_2}, \dotsc, \underbrace{s_m, \dotsc, s_m}_{k_m}), \qquad s_1 > s_2 > \dotsb > s_m,
\end{gather}
and let $K_i = k_1 + k_2 + \dotsb + k_i$, then $\ell(X) = m$, $v_i(X) = s_i$, $n_i(X) = k_i$, and $N_i(X) = K_i$. For~$X$ expressed in~\eqref{eq:general_X}, the master equation~\eqref{eq:abstract_master_eq} can be expressed concretely as
\begin{gather} \label{eq:master_eq_concrete}
 \frac{d}{dt} P(X; t) = \sum^m_{i = 1} \big(1 - q^{p_{s_i - 1}(X^{K_i, -})}\big) a_{s_i - 1} P\big(X^{K_i, -}; t\big) - \sum^m_{i = 1} \big(1 - q^{k_i}\big) a_{s_i} P(X; t),
\end{gather}
where $p_{s_i - 1}(X^{K_i, -})$ is the number of particles on site $s_i - 1$ for the state $X^{K_i, -}$ (as def\/ined in~\eqref{eq:various_notations}). Here we remark that if $X \in \Weylchamber^n$ and $k \in \{ 1, \dotsc, n \}$, $X^{k, -}$ may not be in $\Weylchamber^n$. But all~$X^{K_i, -}$ on the right-hand side of~\eqref{eq:master_eq_concrete} are in~$\Weylchamber^n$.

\subsection{Main result}

In this paper we compute $P_Y(X; t)$, the probability that the particles are in state $X \!= \! (x_1, \dotsc, x_n)$ at time $t > 0$, given the initial condition that the particles are in state $Y = (y_1, \dotsc, y_n)$ at time~$0$. Equivalently, $P_Y(X; t)$ is the solution to the master equation \eqref{eq:abstract_master_eq} with initial condition $P_Y(X; 0) = 1$ if $X = Y$ and $P_Y(X; t) = 0$ if $X \in \Weylchamber^n {\setminus} \{ Y \}$.

To state the main theorem, we introduce some notations. First, for~$X \in \Weylchamber^n$ in the form of~\eqref{eq:general_X}, denote
\begin{gather*}
 W(X) = \prod^{\ell(X)}_{i = 1} \prod^{n_i(X)}_{j = 1} \frac{1 - q^j}{1 - q} = \prod^m_{i = 1} [k_i]_q!,
\end{gather*}
where we use the $q$-deformed factorial $[k]_q! = (1 - q)(1 - q^2) \dotsb (1 - q^k)/(1 - q)^k$. Next, let $w_1, \dotsc, w_n$ be variables. For each pair of integers $\alpha < \beta$, denote
\begin{gather*}
 S_{(\beta, \alpha)} = -\frac{q w_{\beta} - w_{\alpha}}{q w_{\alpha} - w_{\beta}}.
\end{gather*}
For a permutation $\sigma \in S_n$, we say $(\beta, \alpha)$ is an \emph{inversion} of $\sigma$ if $\alpha < \beta$ and $\sigma^{-1}(\alpha) > \sigma^{-1}(\beta)$. Then let
\begin{gather*}
 A_{\sigma}(w_1, \dotsc, w_n) = \prod_{(\beta, \alpha) \text{ is an inversion of } \sigma} S_{(\beta, \alpha)}.
\end{gather*}
We write $A_{\sigma} = A_{\sigma}(w_1, \dotsc, w_n)$ if there is no possibility of confusion, but later we occasionally use dif\/ferent variables for $A_{\sigma}$. Note that for $\sigma = \id \in S_n$, since the identity permutation has no inversion, we let $A_{\id} = 1$. For $X = (x_1, \dotsc, x_n) \in \intZ^n$, $Y = (y_1, \dotsc, y_n) \in \intZ^n$ and $\sigma \in S_n$, def\/ine
\begin{gather} \label{eq:integral_Lambda}
 \Lambda_Y(X; t; \sigma) = \left( \prod^n_{k = 1} \frac{-1}{b_{x_k}} \right) \dashint_{C_1} dw_1 \dotsi \dashint_{C_n} dw_n A_{\sigma} \prod^n_{j = 1} \left[ \prodprime^{x_j}_{k = y_{\sigma(j)}} \left( \frac{b_k}{b_k - w_{\sigma(j)}} \right) e^{-w_j t} \right],
\end{gather}
where the notation $\dashint_{C_i} dw_i$ is a shorthand for $(2\pi i)^{-1} \oint_{C_i} dw_i$ and the notation $\prod'$ is an extension of the usual $\prod$ notation such that
\begin{gather*}
 \prodprime^n_{k = m} f(k) =
 \begin{cases}
 \displaystyle \prod^n_{k = m} f(k) & \text{if $n \geq m$}, \\
 1 & \text{if $n = m - 1$}, \\
\displaystyle \prod^{m - 1}_{k = n + 1} \frac{1}{f(k)} & \text{if $n < m - 1$},
 \end{cases}
\end{gather*}
for example, %$\prod^{\prime -3}_{\phantom{\prime}k = 0}$
 $\prod\limits^{-3}_{k = 0}\!\!{\vphantom{\prod}}' f(k) = 1/[f(-1) f(-2)]$. Later we use the identity
\begin{gather*}
 \left( \prodprime^x_{k = m + 1} f(k) \right) \left( \prodprime^m_{k = y} f(k) \right) = \prodprime^x_{k = y} f(k).
\end{gather*}
With respect to any $w_j$, the integrand in \eqref{eq:integral_Lambda} has explicit poles in the form of $w_j = b_k$ for some~$k$, and the factor $A_{\sigma}$ may introduce poles in the form of $w_j = q w_i$ and $w_j = q^{-1} w_l$ for some~$i$ and~$l$ such that $i < j$ and $l > j$. We require that poles $b_k$ and $q w_i$ are enclosed by~$C_j$, but~$q^{-1} w_l$ are not. To be concrete, we can take that $C_j = \{ \lvert z \rvert = R \}$ for all~$j$ with a large enough~$R$, and choose the counter-clockwise orientation.
\begin{Theorem} \label{thm:main}
 Let $X, Y \in \Weylchamber^n$. Then
 \begin{gather} \label{eq:main_result}
 P_Y(X; t) = \frac{1}{W(X)} \sum_{\sigma \in S_n} \Lambda_Y(X; t; \sigma).
 \end{gather}
\end{Theorem}
